\documentclass[10pt]{article}

% =========================================================
% TEMPLATE
% =========================================================

\input{~/Dropbox/Latex/TemplateExam/template-exam.tex}

% =========================================================
% INFORMACIÓN DEL EXAMEN
% =========================================================

\renewcommand{\CourseName}{Robótica Industrial}
\renewcommand{\CourseCode}{ING 01339 / SOL-ING 220}
\renewcommand{\Institution}{Politécnico Colombiano Jaime Isaza Cadavid}
\renewcommand{\Faculty}{Facultad de Ingeniería}
\renewcommand{\Instructor}{Deimer Miranda Montoya, MSc.(c).}

\renewcommand{\ExamTitle}{Parcial \#1}
\renewcommand{\ExamTopics}{
    Rotaciones 2D y 3D, composición de orientaciones y ejes fijos/móviles
}
\renewcommand{\AcademicPeriod}{2026-2}
\renewcommand{\ExamDuration}{120 minutos}
\renewcommand{\ExamValue}{100 puntos}
\renewcommand{\ExamQuestions}{10}
\renewcommand{\InstitutionLogo}{Template/LogoPoli.png}

%\ExamStandard
\ExamCompact

% =========================================================
\begin{document}

    \pagestyle{empty}

    % =========================================================
    % ENCABEZADO
    % =========================================================

    \ExamHeader
    \StudentInfo

    \ExamInstructions{
        Se permite el uso de una hoja de fórmulas.
        Seleccione una única respuesta en cada pregunta.
        El valor de cada pregunta se indica en el enunciado.
        En las preguntas que involucren cálculo, el procedimiento podrá utilizarse
        como criterio de revisión. Utilice la regla de la mano derecha para establecer
        el sentido positivo de las rotaciones y tenga especial cuidado con el orden
        de multiplicación en las composiciones.
    }

    % =========================================================
    % PREGUNTAS
    % =========================================================

    \begin{multicols}{2}
        \raggedcolumns

        % ---------------------------------------------------------
        % PREGUNTA 1
        % ---------------------------------------------------------

        \begin{question}
            \textbf{[8 puntos]}

            Un vector inicialmente coincide con el eje \(x\) positivo.
            Primero se rota \(35^\circ\) en sentido antihorario y posteriormente
            \(95^\circ\) en sentido horario, ambas rotaciones alrededor del
            mismo origen.

            La rotación equivalente y el vector final son:

            \begin{inlineanswers}
                \answer{
                    \(R(60^\circ)\) y
                    \(\begin{bmatrix}
                        \frac12\\[1mm]
                        \frac{\sqrt3}{2}
                    \end{bmatrix}\)
                }

                \answer{
                    \(R(-60^\circ)\) y
                    \(\begin{bmatrix}
                        \frac12\\[1mm]
                        -\frac{\sqrt3}{2}
                    \end{bmatrix}\)
                }

                \answer{
                    \(R(-130^\circ)\) y
                    \(\begin{bmatrix}
                        -\frac{\sqrt3}{2}\\[1mm]
                        -\frac12
                    \end{bmatrix}\)
                }

                \answer{
                    \(R(130^\circ)\) y
                    \(\begin{bmatrix}
                        -\frac12\\[1mm]
                        \frac{\sqrt3}{2}
                    \end{bmatrix}\)
                }
            \end{inlineanswers}
        \end{question}

        % RESPUESTA DOCENTE: B
        % Procedimiento:
        % theta_eq = 35° - 95° = -60°.
        %
        % Como p = e_x:
        % p_f = R(-60°) e_x
        %     = [cos(-60°), sin(-60°)]^T
        %     = [1/2, -sqrt(3)/2]^T.


        % ---------------------------------------------------------
        % PREGUNTA 2
        % ---------------------------------------------------------

        \begin{question}
            \textbf{[9 puntos]}

            Considere la matriz:

            \[
            R=
            \begin{bmatrix}
                1 & 0 & 0\\
                0 & 0 & -1\\
                0 & 1 & 0
            \end{bmatrix}.
            \]

            ¿Cuál interpretación es correcta?

            \begin{answers}
                \answer{
                    Es una rotación de \(+90^\circ\) alrededor del eje \(x\);
                    la coordenada \(x\) permanece invariante.
                }

                \answer{
                    Es una rotación de \(-90^\circ\) alrededor del eje \(x\);
                    la coordenada \(z\) permanece invariante.
                }

                \answer{
                    Es una rotación de \(+90^\circ\) alrededor del eje \(y\);
                    la coordenada \(y\) permanece invariante.
                }

                \answer{
                    Es una rotación de \(+90^\circ\) alrededor del eje \(z\);
                    la coordenada \(z\) permanece invariante.
                }
            \end{answers}
        \end{question}

        % RESPUESTA DOCENTE: A
        % La matriz coincide con:
        %
        % Rx(90°) =
        % [1  0  0
        %  0  0 -1
        %  0  1  0].
        %
        % El eje x permanece invariante y la rotación ocurre en el plano yz.


        % ---------------------------------------------------------
        % PREGUNTA 3
        % ---------------------------------------------------------

        \begin{question}
            \textbf{[10 puntos]}

            El punto

            \[
            \mathbf{p}
            =
            \begin{bmatrix}
                2\\
                1\\
                3
            \end{bmatrix}
            \]

            se rota \(-90^\circ\) alrededor del eje \(y\).

            El vector resultante es:

            \begin{inlineanswers}
                \answer{
                    \(\begin{bmatrix}
                        -3\\
                        1\\
                        2
                    \end{bmatrix}\)
                }

                \answer{
                    \(\begin{bmatrix}
                        3\\
                        1\\
                        -2
                    \end{bmatrix}\)
                }

                \answer{
                    \(\begin{bmatrix}
                        2\\
                        -3\\
                        1
                    \end{bmatrix}\)
                }

                \answer{
                    \(\begin{bmatrix}
                        -3\\
                        -1\\
                        2
                    \end{bmatrix}\)
                }
            \end{inlineanswers}
        \end{question}

        % RESPUESTA DOCENTE: A
        %
        % Ry(theta) =
        % [ cos(theta)  0  sin(theta)
        %       0       1      0
        %  -sin(theta)  0  cos(theta)]
        %
        % Para theta=-90°:
        %
        % Ry(-90°) =
        % [0 0 -1
        %  0 1  0
        %  1 0  0]
        %
        % Entonces:
        % p' = [-3,1,2]^T.
        %
        % La componente y permanece invariante.


        % ---------------------------------------------------------
        % PREGUNTA 4
        % ---------------------------------------------------------

        \begin{question}
            \textbf{[10 puntos]}

            Un estudiante propone:

            \[
            M=
            \begin{bmatrix}
                1 & 0 & 0\\
                0 & 1 & 0\\
                0 & 0 & -1
            \end{bmatrix}
            \]

            como matriz de rotación.
            \columnbreak
            ¿Cuál afirmación es correcta?

            \begin{answers}
                \answer{
                    \(M\in SO(3)\), porque \(M^TM=I\).
                }

                \answer{
                    \(M\notin SO(3)\), porque no es una matriz ortogonal.
                }

                \answer{
                    \(M\notin SO(3)\): aunque \(M^TM=I\),
                    se cumple \(\det(M)=-1\).
                }

                \answer{
                    \(M\in SO(3)\), porque su inversa coincide con su transpuesta.
                }
            \end{answers}
        \end{question}

        % RESPUESTA DOCENTE: C
        %
        % M^T M = I, por tanto M es ortogonal.
        %
        % Sin embargo:
        % det(M) = -1.
        %
        % Para pertenecer a SO(3) se requiere simultáneamente:
        % R^T R = I
        % det(R) = +1.
        %
        % M representa una reflexión, no una rotación propia.


        % ---------------------------------------------------------
        % PREGUNTA 5
        % ---------------------------------------------------------

        \begin{question}
            \textbf{[10 puntos]}

            La orientación del marco \(\{B\}\) respecto al marco \(\{A\}\) es:

            \[
            {}^A R_B=
            \begin{bmatrix}
                0 & -1 & 0\\
                1 &  0 & 0\\
                0 &  0 & 1
            \end{bmatrix}.
            \]

            ¿Cuál interpretación de sus columnas es correcta?

            \begin{answers}
                \answer{
                    \(x_B=-y_A,\;
                    y_B=x_A,\;
                    z_B=z_A\).
                }

                \answer{
                    \(x_B=y_A,\;
                    y_B=-x_A,\;
                    z_B=z_A\).
                }

                \answer{
                    \(x_B=x_A,\;
                    y_B=z_A,\;
                    z_B=-y_A\).
                }

                \answer{
                    \(x_B=-x_A,\;
                    y_B=-y_A,\;
                    z_B=z_A\).
                }
            \end{answers}
        \end{question}

        % RESPUESTA DOCENTE: B
        %
        % Las columnas de A_R_B corresponden a los ejes de B expresados en A.
        %
        % Columna 1 = [0,1,0]^T  -> x_B = +y_A
        % Columna 2 = [-1,0,0]^T -> y_B = -x_A
        % Columna 3 = [0,0,1]^T  -> z_B = +z_A
        %
        % Corresponde a una rotación de +90° alrededor de z.




        % ---------------------------------------------------------
        % PREGUNTA 6
        % ---------------------------------------------------------

        \begin{question}
            \textbf{[10 puntos]}

            Los marcos \(\{A\}\) y \(\{B\}\) comparten el mismo origen.
            El marco \(\{B\}\) está rotado \(90^\circ\) alrededor de \(z_A\)
            en sentido positivo respecto a \(\{A\}\).

            Se conoce:

            \[
            {}^A\mathbf{p}
            =
            \begin{bmatrix}
                2\\
                1\\
                3
            \end{bmatrix}.
            \]

            Las coordenadas del mismo punto expresadas en \(\{B\}\) son:

            \begin{inlineanswers}
                \answer{
                    \(\begin{bmatrix}
                        -1\\
                        2\\
                        3
                    \end{bmatrix}\)
                }

                \answer{
                    \(\begin{bmatrix}
                        1\\
                        -2\\
                        3
                    \end{bmatrix}\)
                }

                \answer{
                    \(\begin{bmatrix}
                        2\\
                        1\\
                        -3
                    \end{bmatrix}\)
                }

                \answer{
                    \(\begin{bmatrix}
                        -2\\
                        -1\\
                        3
                    \end{bmatrix}\)
                }
            \end{inlineanswers}
        \end{question}

        % RESPUESTA DOCENTE: B
        %
        % A_R_B = Rz(90°).
        %
        % Para cambiar coordenadas de A hacia B:
        %
        % B_p = (A_R_B)^T A_p
        %     = Rz(-90°) A_p.
        %
        % Rz(-90°)[2,1,3]^T
        % = [1,-2,3]^T.
        %
        % El punto físico no se mueve; cambia únicamente su representación.


        % ---------------------------------------------------------
        % PREGUNTA 7
        % ---------------------------------------------------------

        \begin{question}
            \textbf{[10 puntos]}

            Un vector

            \[
            \mathbf{p}
            =
            \begin{bmatrix}
                0\\
                1\\
                0
            \end{bmatrix}
            \]

            experimenta primero una rotación \(R_x(90^\circ)\) y luego una
            rotación \(R_z(90^\circ)\).
            \columnbreak
            ¿Cuál expresión y resultado son correctos?

            \begin{answers}
                \answer{
                    \[
                    \mathbf{p}'
                    =
                    R_x(90^\circ)R_z(90^\circ)\mathbf{p}
                    =
                    \begin{bmatrix}
                        -1\\0\\0
                    \end{bmatrix}
                    \]
                }

                \answer{
                    \[
                    \mathbf{p}'
                    =
                    R_z(90^\circ)R_x(90^\circ)\mathbf{p}
                    =
                    \begin{bmatrix}
                        0\\0\\1
                    \end{bmatrix}
                    \]
                }

                \answer{
                    \[
                    \mathbf{p}'
                    =
                    R_z(90^\circ)R_x(90^\circ)\mathbf{p}
                    =
                    \begin{bmatrix}
                        -1\\0\\0
                    \end{bmatrix}
                    \]
                }

                \answer{
                    \[
                    \mathbf{p}'
                    =
                    R_x(90^\circ)R_z(90^\circ)\mathbf{p}
                    =
                    \begin{bmatrix}
                        0\\0\\1
                    \end{bmatrix}
                    \]
                }
            \end{answers}
        \end{question}

        % RESPUESTA DOCENTE: B
        %
        % La primera transformación aplicada aparece a la derecha:
        %
        % p' = Rz(90°) Rx(90°) p.
        %
        % Primero:
        % Rx(90°)[0,1,0]^T = [0,0,1]^T.
        %
        % Luego:
        % Rz(90°)[0,0,1]^T = [0,0,1]^T.
        %
        % Resultado: [0,0,1]^T.


        % ---------------------------------------------------------
        % PREGUNTA 8
        % ---------------------------------------------------------

        \begin{question}
            \textbf{[10 puntos]}

            Considere el mismo vector:

            \[
            \mathbf{p}
            =
            \begin{bmatrix}
                0\\
                1\\
                0
            \end{bmatrix}.
            \]

            Se comparan:

            \[
            \mathbf{p}_A
            =
            R_z(90^\circ)R_x(90^\circ)\mathbf{p}
            \]

            y:

            \[
            \mathbf{p}_B
            =
            R_x(90^\circ)R_z(90^\circ)\mathbf{p}.
            \]

            ¿Cuál resultado es correcto?

            \begin{answers}
                \answer{
                    \[
                    \mathbf{p}_A
                    =
                    \mathbf{p}_B
                    =
                    \begin{bmatrix}
                        0\\0\\1
                    \end{bmatrix}
                    \]
                }

                \answer{
                    \[
                    \mathbf{p}_A
                    =
                    \begin{bmatrix}
                        0\\0\\1
                    \end{bmatrix},
                    \qquad
                    \mathbf{p}_B
                    =
                    \begin{bmatrix}
                        -1\\0\\0
                    \end{bmatrix}
                    \]
                }

                \answer{
                    \[
                    \mathbf{p}_A
                    =
                    \begin{bmatrix}
                        -1\\0\\0
                    \end{bmatrix},
                    \qquad
                    \mathbf{p}_B
                    =
                    \begin{bmatrix}
                        0\\0\\1
                    \end{bmatrix}
                    \]
                }

                \answer{
                    Ambos resultados tienen que ser iguales porque las matrices
                    de rotación son ortogonales.
                }
            \end{answers}
        \end{question}
            \columnbreak
        % RESPUESTA DOCENTE: B
        %
        % Secuencia A:
        % Rx primero:
        % [0,1,0] -> [0,0,1]
        % Rz después:
        % [0,0,1] -> [0,0,1].
        %
        % Secuencia B:
        % Rz primero:
        % [0,1,0] -> [-1,0,0]
        % Rx después:
        % [-1,0,0] -> [-1,0,0].
        %
        % Por tanto:
        % Rz Rx != Rx Rz.
        %
        % Se evidencia la no conmutatividad de las rotaciones 3D.


        % ---------------------------------------------------------
        % PREGUNTA 9
        % ---------------------------------------------------------

        \begin{question}
            \textbf{[11 puntos]}

            Un marco parte inicialmente alineado con \(\{0\}\).
            Se realiza primero:

            \[
            R_{z_0}(90^\circ)
            \]

            y posteriormente:

            \[
            R_{x_0}(90^\circ).
            \]

            Ambas rotaciones se realizan alrededor de los
            \textbf{ejes fijos del sistema base}.

            ¿Cuál expresión representa correctamente la orientación final?

            \begin{answers}
                \answer{
                    \[
                    R
                    =
                    R_z(90^\circ)R_x(90^\circ)
                    \]
                }

                \answer{
                    \[
                    R
                    =
                    R_x(90^\circ)R_z(90^\circ)
                    \]
                }

                \answer{
                    \[
                    R
                    =
                    R_z(-90^\circ)R_x(90^\circ)
                    \]
                }

                \answer{
                    \[
                    R
                    =
                    R_x(-90^\circ)R_z(-90^\circ)
                    \]
                }
            \end{answers}
        \end{question}

        % RESPUESTA DOCENTE: B
        %
        % Para rotaciones sucesivas alrededor de ejes FIJOS del marco base,
        % la nueva rotación se premultiplica.
        %
        % Estado inicial:
        % R0 = I.
        %
        % Primer giro:
        % R1 = Rz(90°).
        %
        % Segundo giro alrededor de x0:
        % R2 = Rx(90°) R1
        %    = Rx(90°) Rz(90°).
        %
        % La segunda rotación sigue ocurriendo alrededor de x0,
        % no alrededor del eje x que fue arrastrado por el primer giro.


        % ---------------------------------------------------------
        % PREGUNTA 10
        % ---------------------------------------------------------

        \begin{question}
            \textbf{[12 puntos]}

            Un marco parte inicialmente alineado con \(\{0\}\).
            Primero se realiza:

            \[
            R_{z_0}(90^\circ)
            \]

            y posteriormente una rotación de \(90^\circ\) alrededor del
            eje móvil \(x_1\):

            \[
            R_{x_1}(90^\circ).
            \]

            ¿Cuál afirmación describe correctamente la secuencia?

            \begin{answers}
                \answer{
                    La orientación final es
                    \(R_x(90^\circ)R_z(90^\circ)\)
                    porque toda nueva rotación se premultiplica.
                }

                \answer{
                    La orientación final es
                    \(R_z(90^\circ)R_x(90^\circ)\);
                    después del primer giro \(z_1=z_0\) y durante el segundo
                    giro \(x_2=x_1\).
                }

                \answer{
                    La orientación final es
                    \(R_z(90^\circ)R_x(90^\circ)\),
                    pero el segundo giro se realiza alrededor de \(x_0\).
                }

                \answer{
                    Las rotaciones sobre ejes fijos y móviles producen siempre
                    la misma orientación si utilizan los mismos ángulos.
                }
            \end{answers}
        \end{question}

        % RESPUESTA DOCENTE: B
        %
        % En rotaciones sobre ejes MÓVILES la nueva rotación se posmultiplica:
        %
        % R1 = Rz(90°)
        %
        % R2 = R1 Rx(90°)
        %    = Rz(90°) Rx(90°).
        %
        % Después del primer giro alrededor de z0:
        % z1 = z0.
        %
        % Durante el segundo giro se rota alrededor del eje x1,
        % por lo que:
        % x2 = x1.
        %
        % Esta secuencia no es equivalente a la de ejes fijos de la
        % pregunta anterior.

    \end{multicols}

    % =========================================================
    % HOJA DE RESPUESTAS Y CALIFICACIÓN
    % =========================================================

    \AnswerSheet{10}

    \vspace{2mm}

    {\color{ExamGreen}\hrule height 0.6pt}

    \vspace{2mm}

    \ExamScore

\end{document}